% ============================================================================
% Prospectos Solares. Sistema visual de la presentacion.
% ----------------------------------------------------------------------------
% La paleta, el isotipo y las proporciones se toman del reporte de lotes
% (reporte_predios/plantilla.py, modo claro), para que la presentacion y el
% entregable se lean como un mismo producto y no como dos piezas distintas.
%
% Reglas del sistema:
%   Color      un solo verde azulado de marca. El acento oscuro (acento) es el
%              que lleva texto y linea; el claro (marca) solo identifica.
%              bien, aviso y mal se reservan para estado, nunca para decorar.
%   Rejilla    margen lateral unico (\margenlat) y altura util unica
%              (\altoutil). Ninguna lamina cambia de proporcion.
%   Filete     un solo divisor de marca: 1,8 cm por 1,4 pt en acento. Las
%              separaciones internas usan la hairline de 0,6 pt en linea.
%   Tipografia una sola familia de palo seco, tres pesos de gris (tinta,
%              tinta2, tinta3) y jerarquia por tamano, no por color.
%   Cajas      toda caja de diagrama es un nodo de tikz con text width y
%              minimum height: la caja se dimensiona a partir del texto, de
%              modo que el texto nunca puede desbordarla.
%   Lectura    sin justificar y sin partir palabras.
% ============================================================================

\usetheme{default}
\usefonttheme{professionalfonts}
\setbeamertemplate{navigation symbols}{}

% --- Paleta ------------------------------------------------------------------
\definecolor{marca}{HTML}{18919C}        % identidad
\definecolor{acento}{HTML}{0F6A72}       % texto y linea de acento
\definecolor{acento2}{HTML}{0B5057}      % acento fuerte, enfasis
\definecolor{acentosuave}{HTML}{E3EFF0}  % fondo de acento
\definecolor{tinta}{HTML}{1E2829}        % texto principal
\definecolor{tinta2}{HTML}{49585A}       % texto secundario
\definecolor{tinta3}{HTML}{5C6B6D}       % texto terciario, notas
\definecolor{linea}{HTML}{DCE2E2}        % hairline
\definecolor{fondo}{HTML}{F4F6F6}        % fondo de pagina
\definecolor{superficie}{HTML}{EAEFF0}   % superficie elevada
\definecolor{bien}{HTML}{2F6B4C}         % estado favorable
\definecolor{aviso}{HTML}{8A5C17}        % estado en curso
\definecolor{mal}{HTML}{9C3A1F}          % estado que impide

% Clasificacion, misma escala que el reporte de lotes.
\definecolor{clasei}{HTML}{17563E}
\definecolor{claseii}{HTML}{2C5F9E}
\definecolor{claseiii}{HTML}{C08322}
\definecolor{claseiv}{HTML}{879596}

\setbeamercolor{normal text}{fg=tinta, bg=white}
\setbeamercolor{background canvas}{bg=white}
\setbeamercolor{frametitle}{fg=tinta, bg=white}
\setbeamercolor{framesubtitle}{fg=tinta2}
\setbeamercolor{structure}{fg=acento}
\setbeamercolor{alerted text}{fg=acento2}

\setbeamerfont{frametitle}{size=\large, series=\bfseries}
\setbeamerfont{framesubtitle}{size=\scriptsize, series=\mdseries}

% --- Rejilla fija ------------------------------------------------------------
% Todas las laminas comparten el mismo margen y la misma altura util, para que
% las proporciones no bailen entre una y otra.
\newlength{\margenlat}\setlength{\margenlat}{1.0cm}
\newlength{\altoutil}\setlength{\altoutil}{0.615\paperheight}

% Sin partir palabras: en laminas con columnas estrechas el corte silabico
% ensucia la lectura y en una presentacion no compensa.
\hyphenpenalty=10000
\exhyphenpenalty=10000
\setbeamersize{text margin left=\margenlat, text margin right=\margenlat}

% --- Grosores del sistema ----------------------------------------------------
\newlength{\anchofilete}\setlength{\anchofilete}{1.8cm}
\newlength{\altofilete}\setlength{\altofilete}{1.4pt}
\newlength{\hairline}\setlength{\hairline}{0.6pt}

% Divisor de marca. Es la unica linea gruesa de toda la presentacion.
\newcommand{\reglamarca}[1][acento]{\textcolor{#1}{\rule{\anchofilete}{\altofilete}}}

% Hairline de ancho libre, para separar bloques dentro de una lamina.
\newcommand{\filete}[1][\linewidth]{\textcolor{linea}{\rule{#1}{\hairline}}}

% Bandera moderada. \raggedright deja estiramiento infinito al margen derecho y
% en columna estrecha produce lineas de longitudes muy dispares. Acotando el
% estiramiento a dos cuadratines las lineas se igualan bastante sin justificar,
% que es lo que el sistema no quiere porque abre rios y partiria palabras.
\newcommand{\bandera}{\rightskip=0pt plus 2em\relax\parindent=0pt}

% Rotulo breve sobre un bloque. Siempre en versales y en gris terciario.
\newcommand{\rotulo}[2][tinta3]{{\scriptsize\textcolor{#1}{\MakeUppercase{#2}}}}

% --- Isotipo: siete barras de altura creciente, extremos redondeados ---------
\newcommand{\isotipo}[2][marca]{%
  \begin{tikzpicture}[y=#2, x=#2]
    \foreach \x/\y/\h in {0/0.270/0.460, 0.152/0.215/0.570, 0.304/0.160/0.680,
                          0.457/0.110/0.780, 0.609/0.060/0.880, 0.761/0.025/0.950,
                          0.909/0/1.0}
      {\fill[#1, rounded corners=0.0456*#2] (\x, \y) rectangle (\x+0.0913, \y+\h);}
  \end{tikzpicture}%
}

\newcommand{\marcaTexto}[1][marca]{%
  \raisebox{-0.28em}{\isotipo[#1]{0.62em}}\hspace{0.5em}%
  \textcolor{#1}{\textbf{\footnotesize MÉTODOS MIXTOS}}\,%
  \textcolor{tinta3}{\scriptsize CONSULTORES}%
}

% --- Titulo de lamina --------------------------------------------------------
% Titulo, subtitulo opcional y el filete de marca. Las medidas verticales son
% las mismas en todas las laminas: el cuerpo siempre arranca a la misma altura.
\makeatletter
\setbeamertemplate{frametitle}{%
  \vspace{0.45em}%
  \begin{beamercolorbox}[wd=\paperwidth, leftskip=\margenlat, rightskip=\margenlat]{frametitle}%
    % Sin \raggedright: su definicion pone \leftskip a cero y con el se llevaba
    % la sangria del propio beamercolorbox. Los titulos salian pegados al borde
    % del papel mientras el cuerpo arrancaba en el margen. Se compone en bandera
    % fijando las dos sangrias a mano.
    % El \rightskip se fija con \dimexpr: \margenlat es un registro de glue y
    % detras de un registro TeX no lee el "plus", que se imprimia como texto.
    \leftskip=\margenlat
    \rightskip=\dimexpr\margenlat\relax plus 1fil\relax
    \parindent=0pt%
    \usebeamerfont{frametitle}\insertframetitle%
    \ifx\insertframesubtitle\@empty\else%
      \\[0.25em]{\usebeamerfont{framesubtitle}\usebeamercolor[fg]{framesubtitle}\insertframesubtitle}%
    \fi%
  \end{beamercolorbox}%
  \vspace{0.5em}%
  % Sin \hspace: fuera del beamercolorbox ya rige el margen de lamina, de modo
  % que la sangria se sumaba dos veces y el filete arrancaba a 2 cm mientras el
  % titulo y el cuerpo arrancaban a 1 cm.
  \reglamarca%
  % Aire bajo el filete. Estaba en 0,85 em y dejaba una banda muerta entre el
  % filete y la figura justo donde la figura tenia que empezar, lo que obligaba a
  % encoger despues. Con 0,35 em la lamina recupera alto util y varias figuras
  % bajan de shrink.
  \vspace{0.35em}%
}
\makeatother

% --- Pie ---------------------------------------------------------------------
% Seccion a la izquierda, isotipo y numero a la derecha. Sin linea: el pie se
% separa por aire, no por trazo.
\setbeamertemplate{footline}{%
  \hbox{%
  \begin{beamercolorbox}[wd=0.5\paperwidth, ht=2.2ex, dp=1.6ex, leftskip=\margenlat]{}%
    \tiny\textcolor{tinta3}{\MakeUppercase{\insertsectionhead}}%
  \end{beamercolorbox}%
  \begin{beamercolorbox}[wd=0.5\paperwidth, ht=2.2ex, dp=1.6ex, rightskip=\margenlat]{}%
    \hfill\raisebox{-0.15em}{\isotipo[linea]{0.44em}}\hspace{0.5em}\tiny\textcolor{tinta3}{\insertframenumber}%
  \end{beamercolorbox}}%
}

% Sin justificar: en columnas estrechas la justificacion abre rios de espacio y
% parte palabras. La alineacion se marca dentro de cada minipage. No enganchar
% \frame con \apptocmd: rompe la lectura de [plain] y de los titulos.

\setbeamertemplate{itemize item}{\raisebox{0.14em}{\textcolor{acento}{\rule{0.4em}{0.4em}}}}
\setbeamertemplate{itemize subitem}{\raisebox{0.16em}{\textcolor{tinta3}{\rule{0.32em}{0.32em}}}}
\setlength{\leftmargini}{1.3em}
\setlength{\parskip}{0.4em}

% Tablas: hairline del sistema y aire suficiente entre filas.
\setlength{\heavyrulewidth}{\altofilete}
\setlength{\lightrulewidth}{\hairline}
\setlength{\aboverulesep}{0.5ex}
\setlength{\belowrulesep}{0.7ex}

% ============================================================================
% Piezas
% ============================================================================

% --- Separador de parte ------------------------------------------------------
% Lamina limpia: filete, rotulo de etapa, titulo grande y una linea de encuadre.
%
% Sin [plain]: con el, estas laminas salian sin pie y sin numero, y la numeracion
% visible saltaba de 5 a 7, de 8 a 10 y asi cinco veces. En un documento que se
% cita por lamina y sobre el que se van a mandar comentarios por escrito, eso
% confunde de inmediato. El pie queda, con su numero.
\newcommand{\parte}[3]{%
  \begin{frame}
    \vfill
    \hspace{\margenlat}\begin{minipage}{0.78\paperwidth}
      \raggedright
      \reglamarca\\[1.0em]
      \rotulo{#1}\\[0.35em]
      {\fontsize{26}{30}\selectfont\textcolor{tinta}{\textbf{#2}}}\\[0.7em]
      {\small\textcolor{tinta2}{#3}}
    \end{minipage}
    \vfill
  \end{frame}
}

% --- Lamina de afirmacion ----------------------------------------------------
% Una idea sostenida por un parrafo, con mucho aire alrededor.
\newcommand{\afirmacion}[2]{%
  \vspace{0.3em}
  \begin{minipage}{0.88\textwidth}
    \raggedright
    {\fontsize{15}{19}\selectfont\textcolor{tinta}{#1}}\\[0.9em]
    {\footnotesize\textcolor{tinta2}{#2}}
  \end{minipage}}

% --- Imagen a caja fija ------------------------------------------------------
% Tres cajas y solo tres, en vez de catorce alturas escritas a mano por lamina.
% La figura se dimensiona por la caja, de modo que dos laminas seguidas no
% cambian de peso. keepaspectratio hace que mande la restriccion mas estrecha.
%
%   \imagen        ancha, ocupa el ancho util entero
%   \imagenmedia   media, para la columna mayor de una lamina a dos columnas
%   \imagenalta    alta, para las capturas en retrato, donde manda la altura
\newcommand{\imagen}[1]{%
  \begin{center}\includegraphics[width=\textwidth, height=\altoutil, keepaspectratio]{#1}\end{center}}

\newcommand{\imagenmedia}[1]{%
  \includegraphics[width=\linewidth, height=0.52\paperheight, keepaspectratio]{#1}}

\newcommand{\imagenalta}[1]{%
  \includegraphics[width=\linewidth, height=0.66\paperheight, keepaspectratio]{#1}}

% --- Cifra -------------------------------------------------------------------
% Dato y su leyenda. La hairline superior las alinea cuando van varias en fila.
\newcommand{\cifra}[2]{%
  \begin{minipage}[t]{0.22\textwidth}
    \raggedright
    \filete\\[0.55em]
    {\color{acento}\fontsize{24}{26}\selectfont\textbf{#1}}\\[0.3em]
    {\scriptsize\color{tinta2} #2}
  \end{minipage}}

% Cifra de anclaje: la misma pieza en tamano de titular, para sostener una
% lamina de \afirmacion sin inventar contenido. Lleva la cifra que la frase ya
% menciona, no un dato nuevo.
\newcommand{\cifraancla}[3][0.30\textwidth]{%
  \begin{minipage}[t]{#1}
    \raggedright
    \filete\\[0.6em]
    {\color{acento}\fontsize{22}{25}\selectfont\textbf{#2}}\\[0.45em]
    {\scriptsize\color{tinta2} #3}
  \end{minipage}}

% Rotulo con hairline superior, sin cifra: para enumerar condiciones en fila.
\newcommand{\rotuloancla}[2][0.21\textwidth]{%
  \begin{minipage}[t]{#1}
    \raggedright
    \filete\\[0.6em]
    {\small\color{tinta}\textbf{#2}}
  \end{minipage}}

% --- Chip --------------------------------------------------------------------
% Etiqueta de estado. Nodo de tikz para que la caja se ajuste al texto y quede
% alineada con la linea base.
\newcommand{\chip}[2]{%
  \tikz[baseline=(c.base)]{%
    \node[rectangle, rounded corners=2pt, inner xsep=5pt, inner ysep=2.6pt,
          fill=#1!10, text=#1, font=\scriptsize\bfseries] (c) {#2};}}

% --- Cajas de diagrama -------------------------------------------------------
% El nodo ES la caja: el texto no puede desbordarla porque la caja se
% dimensiona a partir del texto. Dibujar rectangulo y texto por separado
% permite que el texto se salga, que es el defecto que esto corrige.
% El argumento del estilo es el ancho de texto. El inner sep de 7 pt es parte de
% la rejilla de los diagramas: cambiarlo estrecha el hueco de las flechas.
\tikzset{
  caja/.style={
    rectangle, rounded corners=2pt, draw=none,
    text width=#1, align=flush left, inner sep=7pt,
    anchor=north west, minimum height=3.55cm},
  cajarot/.style={
    rectangle, rounded corners=2pt, draw=linea, line width=\hairline,
    text width=#1, align=flush left, inner sep=7pt,
    anchor=north west, minimum height=3.55cm},
}

% Contenido tipico de una caja: rotulo, titulo y descripcion.
\newcommand{\cajatexto}[3]{%
  {\scriptsize\textcolor{acento}{\textbf{\MakeUppercase{#1}}}}\\[0.35em]
  {\small\textcolor{tinta}{\textbf{#2}}}\\[0.35em]
  {\scriptsize\textcolor{tinta2}{#3}}}

% --- Nota de fuente ----------------------------------------------------------
% Siempre al pie del contenido, en gris terciario y sin justificar.
% A 7 pt y no en scriptsize: la cita de fuente competia con el cuerpo de la
% lamina cuando iba debajo de una imagen. Es el minimo del sistema, el mismo del
% pie de lamina, de modo que sigue siendo legible proyectada.
\newcommand{\fuente}[1]{\vspace{0.3em}\par{\raggedright\fontsize{7}{8.4}\selectfont\color{tinta3} #1\par}}
